\chapter{关键代码索引与复现命令}

\section{关键代码索引}
表 \ref{tab:appendix-code-index} 给出本报告最常引用的代码段，便于读者在仓库中快速定位。

\begin{longtable}{p{0.26\textwidth} p{0.18\textwidth} p{0.46\textwidth}}
\caption{关键代码索引} \label{tab:appendix-code-index} \\
\toprule
文件 & 行号 & 内容 \\
\midrule
\endfirsthead
\toprule
文件 & 行号 & 内容 \\
\midrule
\endhead
\code{SPH\_Poiseuille.m} & 16--29 & 路径初始化与环境变量覆盖输出接口 \\
\code{SPH\_Poiseuille.m} & 32--36 & 自动编译两个 MEX 模块 \\
\code{SPH\_Poiseuille.m} & 38--84 & 读取配置与派生参数 \\
\code{SPH\_Poiseuille.m} & 86--117 & 流体粒子和单层 shell 壁面初始化 \\
\code{SPH\_Poiseuille.m} & 119--155 & restart 检查与恢复 \\
\code{SPH\_Poiseuille.m} & 157--234 & 初始邻居搜索、密度修正与运行结构体组织 \\
\code{SPH\_Poiseuille.m} & 236--294 & 主循环整体结构 \\
\code{SPH\_Poiseuille.m} & 244--255 & 统一时间步与高层 MEX 门面调用 \\
\code{SPH\_Poiseuille.m} & 257--273 & 周期包裹、壁面防穿透、排序和邻居重建 \\
\code{SPH\_Poiseuille.m} & 286--293 & restart 保存与中截面剖面记录 \\
\code{SPH\_Poiseuille.m} & 296--506 & 后处理、误差评估和绘图 \\
\code{SPH\_Poiseuille.m} & 524--549 & 邻居缓存与高层 MEX 桥接函数 \\
\code{SPH\_Poiseuille.m} & 637--643 & 统一 CFL 时间步函数 \code{cfl\_time\_step} \\
\code{SPH\_Poiseuille.m} & 687--704 & 周期边界函数 \code{periodic\_bounding} \\
\code{SPH\_Poiseuille.m} & 739--757 & 壁面防穿透函数 \code{bounding\_from\_wall} \\
\code{build\_shell\_wall\_particles.m} & 10--19 & 单层 shell 壁面几何生成 \\
\code{mex/sph\_neighbor\_search\_mex.c} & 43--60 & cubic spline 核函数与径向导数 \\
\code{mex/sph\_neighbor\_search\_mex.c} & 112--337 & cell-linked list 邻居搜索与最小像距 \\
\code{mex/sph\_physics\_mex.c} & 89--362 & 密度重初始化、体积与核梯度修正 \\
\code{mex/sph\_physics\_mex.c} & 384--538 & 粘性力与单层 shell 无滑移接触 \\
\code{mex/sph\_physics\_mex.c} & 556--696 & transport correction \\
\code{mex/sph\_physics\_mex.c} & 717--947 & 第一阶段积分 \\
\code{mex/sph\_physics\_mex.c} & 964--1097 & 第二阶段积分 \\
\code{mex/sph\_physics\_mex.c} & 1116--1326 & 高层推进门面 \code{advance\_shell\_step} \\
\code{mex/sph\_physics\_mex.c} & 1339--1429 & 壁面剪应力监测门面 \code{wall\_shear\_monitor} \\
\code{mex/sph\_physics\_mex.c} & 1431--1454 & mode 分发入口 \code{mexFunction} \\
\code{config.ini} & 4--19 & 物理参数与仿真控制参数 \\
\bottomrule
\end{longtable}

\section{推荐运行命令}
项目当前以 MATLAB 为主驱动，推荐命令如下。

\begin{lstlisting}[style=codeblock,language=bash,caption={主算例运行命令}]
matlab -batch "run('SPH_Poiseuille.m')"
\end{lstlisting}

\begin{lstlisting}[style=codeblock,language=bash,caption={轻量 smoke test}]
matlab -batch "run('tests/test_poiseuille_smoke.m')"
\end{lstlisting}

\begin{lstlisting}[style=codeblock,language=bash,caption={MEX mode 接口测试}]
matlab -batch "run('tests/test_sph_physics_modes.m')"
\end{lstlisting}

\begin{lstlisting}[style=codeblock,language=bash,caption={技术报告编译命令}]
cd docu
latexmk -xelatex -interaction=nonstopmode -file-line-error sph_poiseuille_report.tex
\end{lstlisting}

\section{附录说明}
这份附录只列出最关键的代码与命令。如果后续需要继续扩展到性能 profiling、参数扫描或不同边界组合比较，建议先在仓库中把运行输出结构化成单独的数据目录，再在报告附录中增加对应的数据文件索引。
